\documentclass[10pt]{article}
\usepackage[margin=20mm]{geometry}
\usepackage{amsmath,amssymb,mathtools,fontspec,unicode-math,xurl,hyperref}
\IfFontExistsTF{DejaVu Serif}{\setmainfont{DejaVu Serif}}{\setmainfont{Latin Modern Roman}}
\setmathfont{latinmodern-math.otf}
\setlength{\parindent}{0pt}\setlength{\parskip}{6pt}\sloppy
\title{Initial scientific review: M25-0036}
\author{}
\date{25 October 2026}
\begin{document}
\maketitle
\section*{Identity, version, and outcome}
Alexander Smith, The distribution of \(\ell^\infty\)-Selmer groups in degree \(\ell\) twist families II, DOI 10.1090/jams/1063. The reviewed object is arXiv:2207.05143v2, submitted 8 February 2023, internal date 9 February 2023, 65 pages, together with the supplied author TeX. PDF SHA-256:\par
> d8f2b5f5d033d7163d292da86d794f4a3e2
> b780755c5eebb6f3ce093a68abcee.\par
Page numbers are one-based PDF pages and agree with printed numbers. There was no journal-final comparison.\par
Overall judgment: Major for the false interpolation estimate used in the proof of the main theorem, with three local Minor findings. The main Selmer distribution theorem is not disproved. A weaker interpolation bound is proved below and is compatible with the relevant asymptotic scales. Severity refers to the demonstrated invalid core derivation in the supplied version, irrespective of repair difficulty.\par
Text pages 1–23 and 56–62 were read (30 pages). Original images of pages 1–3, 6–22, and 56–62 were inspected (27 pages). Selected author source was inspected directly: `introB.tex` lines 16–40, `results.tex` lines 121–178, and `distribution.tex` lines 93–129, with additional targeted search context. This is a claim-focused review, not certification of the full analytic machinery.\par
\section*{Main claims and dependency map}
Theorem 2.14 compares the empirical distribution of entire nonincreasing Selmer-rank sequences with a Markov law in total variation. It requires admissible twists, fixed globally realizable local data, favored twists, and either Case 2.12 or Case 2.13. Conditions on cofavored submodules, connecting-map automorphisms, universal Selmer elements, and parity cannot be omitted. For higher prime-power twisting groups, admissibility requires full inertia at every new ramified prime.\par
The core route is cofavored-module classification (Section 4), local character expressions and negligible terms (Sections 5–7), gridding and rough moments (Sections 8–9), refined moments (Proposition 10.3), moment identification (10.4), coefficient recovery (10.5), and Part I’s higher-grid distribution theorem. Theorem 2.14’s proof explicitly invokes Proposition 10.5 on page 61. This is the dependency supporting the Major finding.\par
The average-rank application in Theorem 1.1 instead uses Theorem 2.6 and the isogeny construction of Example 3.9. Its small-parameter defect below does not establish failure of bounded average rank.\par
\section*{Problem 1 (Major): Proposition 10.5 has false uniform coefficient bounds}
\subsection*{Claim, location, and failed step}
On pages 59–60, Proposition 10.5 asserts a universal \(C>0\). For real coefficients \(a_i\), an integer \(m\ge1\), a rational prime \(\ell\), and finite \[B=\sum_{i\ge0}|a_i|\ell^{mi},\] small values of \(f(z)=\sum_i a_i z^i\) on \(1,\ell,\ldots,\ell^{m-1}\) are claimed to imply \[|a_i|\le \ell^{(-i+C)m}(B\ell^{-m^2/2}+\epsilon)\quad(i\ge0).
 \tag{R1}\] With both signs of those nodes, the gain is claimed to be \(\ell^{-m^2}\) instead. The exact statements appear in `distribution.tex` lines 105 and 107.\par
The proof correctly obtains a Cauchy estimate for coefficients of \(q=(f-g)/p\), where \(p\) is the node polynomial and \(g\) is the interpolation polynomial. It then transfers the factor \(\ell^{-im}\) to coefficients of \(f=g+pq\) merely because \(p\) has bounded coefficients (source line 125). This transfer is invalid: the coefficient convolution includes \(b_{i-r}\), whose bound has the additional factor \(\ell^{rm}\).\par
\subsection*{Counterexample to the first bound, for every proposed constant}
Fix \(\ell=2\) and take \[f_m(z)=\prod_{r=0}^{m-1}(1-z/2^r).\] This polynomial has real coefficients, finite \(B\), and is zero at every required node. Its leading coefficient and weighted coefficient sum are \[|a_m|=2^{-m(m-1)/2},\qquad
 B=\prod_{s=1}^{m}(1+2^s)
 =2^{m(m+1)/2}\prod_{s=1}^{m}(1+2^{-s})
 <e\,2^{m(m+1)/2}.\] For \(i=m\), the ratio of the \(B\) term on the right of (R1) to \(|a_m|\) is less than \[e\,2^{-m^2/2+Cm}.\] For any fixed proposed \(C\), choose an integer \(m\) sufficiently large that this is less than \(1/2\). Choose \[0<\epsilon\le
 \min\{1,\ |a_m|/[4\,2^{(-m+C)m}]\}.\] All node conditions still hold, while the full claimed right side is less than \(3|a_m|/4\). Thus no universal \(C\) can make (R1) true. This verifies all hypotheses rather than relying on one guessed value of the unspecified constant.\par
\subsection*{Counterexample to the signed-node bound}
Use \[f_m(z)=\prod_{r=0}^{m-1}(1-z^2/2^{2r}).\] It vanishes at all \(2m\) signed nodes, has \[|a_{2m}|=2^{-m(m-1)},\qquad
 B=\prod_{s=1}^{m}(1+2^{2s})
 <e^{1/3}2^{m(m+1)}.\] For \(i=2m\), the ratio of the claimed \(B\) term to the actual coefficient is at most \(e^{1/3}2^{-m^2+Cm}\), again tending to zero. The same positive-\(\epsilon\) choice finishes the contradiction. Exact rational computations for \(m=8\), \(\ell=2\), illustrative \(C=1\), and \(\epsilon=1\) are saved for both parts; the analytic families, not those finite examples alone, disprove universality.\par
\subsection*{Counterevidence: a weaker sufficient-range estimate}
A valid replacement, for \(0\le i\le m\), is \[|a_i|\le \ell^{C'm}(B\ell^{-m^2/2}+\epsilon)
 \tag{R2}\] in the first case, and the same expression with \(\ell^{-m^2}\) in the signed case, for an absolute \(C'\).\par
Here is a proof of the stated replacement. The manuscript’s interpolation construction gives, on the disk of radius \(\ell^m\), \[|q|\le B\ell^{-m(m-1)/2}
       +4m\epsilon\ell^{2m+1}.\] On \(|z|=1/2\), the node polynomial has magnitude at most \(\prod_{r\ge0}(1+(2\ell^r)^{-1})\le e\). The Lagrange denominators have magnitude at least \(1/4\), so \(|g|\le4em\epsilon\) there. Hence \(|f|\le e|q|+4em\epsilon\) on this smaller circle. Cauchy’s estimate contributes \(2^i\le\ell^m\) for \(i\le m\), and all remaining factors are absorbed in \(\ell^{C'm}\), proving (R2).\par
For signed nodes write \(f(z)=F(z^2)+zG(z^2)\). Both \(F\) and \(G\) are bounded by \(\epsilon\) at \(\ell^{2r}\), because \(F(\ell^{2r})=(f(\ell^r)+f(-\ell^r))/2\) and \(G(\ell^{2r})=(f(\ell^r)-f(-\ell^r))/(2\ell^r)\). Their weighted coefficient sums for base \(\ell^2\) are at most \(B\) and \(B/\ell^m\). Applying the preceding proof, which works for every real base at least two, gives the signed version of (R2).\par
On page 61 the relevant \(B\) grows like \(\ell^{m^2/2+O(m)}\) in the alternating case and \(\ell^{m^2/4+O(m)}\) otherwise. The signed or unsigned replacement still gives exponential decay in \(m^2\). The target coefficient index is at most \((\log^{(3)}H)^{1/4}\), far smaller than an admissible interpolation order of size \((\log\log H)^{1/8-\eta}\). Thus this local repair is compatible with the required range and scale, after choosing margins in \(\eta\). This does not certify the entire remaining arithmetic proof or all constants.\par
The supported conclusion is that an explicitly used core lemma and its final coefficient step are false as written. It is not a counterexample to Theorem 2.14 or to the arithmetic rank distributions.\par
\section*{Problems 2–4: separate local defects}
\subsection*{Problem 2 (Minor): finite alternating transition formula}
Case 2.13, page 9, and `results.tex` line 161 display a numerator \(\prod_{k=1}^{j}(1-2^{-n-j+k})\). For \(n=j=2\), the two alternating matrices are the zero matrix and \(\left(\begin{smallmatrix}0&1\\1&0\end{smallmatrix}\right)\). Therefore \(P(2\mid2)=P(0\mid2)=1/2\). The printed formula instead gives \(P(2\mid2)=7/8\), so its row sums to \(11/8\).\par
Counting the radical and then a nondegenerate alternating form on the quotient gives the corrected formula \[P(j\mid n)=2^{-j(j-1)/2}
 \prod_{k=1}^{j}\frac{1-2^{-n+j-k}}{1-2^{-k}}
 \prod_{k=1}^{(n-j)/2}(1-2^{-2k+1})
 \tag{R3}\] for \(n-j\) even, and zero otherwise. This follows from the Gaussian binomial count of \(j\)-subspaces and the number of nondegenerate alternating forms; the power of two simplifies to \(-j(j-1)/2\). Exact enumeration for every \(n\le5\) agrees with (R3).\par
Counterevidence is substantial: the preceding verbal definition uses genuine uniform random matrices, Part I’s Figure 1.1 gives the correct finite probabilities, and the fixed-\(j\) limiting numerator tends to one in either expression. Proposition 10.4 counts kernels directly. This is a local erroneous formula, with a verified replacement, rather than evidence of an incorrect intended Markov model. Any use of the displayed finite formula requires correction.\par
\subsection*{Problem 3 (Minor): the small-positive-parameter range of Theorem 1.1}
Page 2 and `introB.tex` lines 20–21 quantify over every \(m>0\) without an integer restriction and bound the mean of \(e^{m\operatorname{rank}(A/K)}\) by \(e^{Cm^2}\). This cannot hold for all positive real \(m\).\par
Take \(F=\mathbb Q\), \(\ell=2\), and \(E:y^2=x^3-2\). The point \((3,5)\) is nontorsion: good-reduction counts at \(5\) and \(7\) are \(6\) and \(7\), which, by prime-to-residue-characteristic injectivity, force rational torsion to be trivial. The curve is nonsingular, with discriminant \(-1728\). Thus \(\operatorname{rank}(E/K)\ge1\) for every quadratic extension \(K\). For any proposed \(c,C>0\), choose \(H>\max(C,8,e^e)\), so the extension set is nonempty, and choose \(0<m<\min(1/C,c\log\log\log H)\). Then its mean is at least \(e^m>e^{Cm^2}\).\par
Theorem 2.6 explicitly uses a positive integer moment order, and Example 3.9 applies it to obtain boundedness. That is the natural intended reading. Restricting to \(m\ge1\), or explicitly to positive integer \(m\), removes this counterexample. If a real-parameter bound for all \(m>0\) is desired, a linear term in the exponent is needed: a bound at \(m=1\) gives \(e^{Dm}\) for \(0<m\le1\) by concavity, and together with a large-\(m\) estimate gives \(e^{D(m+m^2)}\). No failure of the bounded-average-rank conclusion is inferred. The literal endpoint quantifier, not the existence of a useful moment bound, is the defect.\par
\subsection*{Problem 4 (Minor): favored density needs the subsequent hypotheses}
Immediately after Example 1.2, page 3, the density of favored squarefree \(d\) is stated to be \(1/2\) under only \(b(a^2-4b)\ne0\). Take \(a=0,b=-1\). Then \(E:y^2=x^3-x\) is nonsingular, \(K=\mathbb Q\), and \(K_0=\mathbb Q(i)\). At every eligible prime, splitting in \(K_0\) contributes at most the splitting count in the trivial extension \(K\). Every squarefree \(d\) is favored, giving density one.\par
The following Proposition 1.3 and Theorem 1.4 explicitly require \(K,K_0\) to be distinct nontrivial quadratic extensions. The example violates those later hypotheses and does not refute either result. The earlier density sentence should carry those restrictions; the repair of the full density theorem is not reproved here.\par
\section*{Positive checks and qualifications}
\subsection*{Favored twists and cofavored modules}
The exact-sequence cardinality calculation in Proposition 2.5, page 18, gives the displayed local Tamagawa factor after new ramification forces invariants into \(\omega\)-torsion. The global lower bound remains conditional on the cited Euler-characteristic theorem.\par
For Proposition 4.3, positivity of the common weight makes equality of the summed invariant-dimension inequalities force equality at each group element. Adding a positive constant to a superlative is legitimate using the unweighted inequality. For Proposition 4.5, a submodule with zero intersection with the second summand and full first projection is a graph. Its connecting map is the difference \(\beta\delta_1-\delta_2\beta\), giving precisely the commuting condition. The resulting finite division-ring argument and the subspace description explain the Gaussian-binomial contribution to the moments. The induction in Proposition 4.7 must be read on \(\omega\)-torsion; local omitted torsion notation and the reversed sign in the final dimension display are not treated as separate substantive findings. With the correct sign, varying the quotient dimension gives the stated contradiction.\par
\subsection*{Moments, parity, and tails}
For \(m\) vectors spanning dimension \(d\), an alternating random matrix kills them with probability \(2^{-nd+d(d+1)/2}\). Counting the span relations gives \[\operatorname{gr}_2(d,m)\prod_{r=0}^{d-1}(2^n-2^r)\;\text{tuples}.\] This validates the finite counting step of Proposition 10.4 and gives limiting first moments \(3\) and \(15\) for orders one and two. In the rectangular case the first moment is \(1+\ell^u\).\par
Parity information is essential: equal unsigned moments alone do not select between the even and odd limiting laws. The signed-node step on page 61 uses fixed parity. For an odd-degree hyperelliptic polynomial with full symmetric Galois action, the fixed-space dimension on \(\mathbb F_2^{2g+1}/\langle(1,\ldots,1)\rangle\) is the number of permutation cycles minus one. Its parity is the sign character, consistent with the parity-invariant designation in Example 3.7. This does not independently verify the cited rational-point bounds in Theorem 3.8.\par
The rough moment bound yields a rank tail of size \(\exp(-cR^2)\) by choosing an integer moment proportional to \(R\). At \(R=(\log^{(3)}H)^{1/4}\) this moment lies within Theorem 2.6’s permitted logarithmic range and gives the required \(\exp(-c\sqrt{\log^{(3)}H})\) scale, conditional on the positive density of the fixed-local favored family.\par
Part I’s rational higher-torsion issue is recorded separately in M25-0035. Part II explicitly imposes new ramification in Remark 2.11; its admissible nontrivially ramified twists force rational invariants to be killed by \(\omega\). The finite exceptional family does not itself damage an asymptotic theorem. The class-group roots-of-unity exclusion and \(u=-r_2-r'_1\) agree with Part I. The precise global duality, parity, Chebotarev, and gridding estimates have not all been independently verified.\par
\section*{Impact, repairs, and remaining work}
Overall | Major for Proposition 10.5’s directly used invalid estimate; three separate local Minor findings.\par
Core conclusion | Theorem 2.14 not disproved; a weaker coefficient estimate compatible with its scales is proved here.\par
Repair status | (R2) and (R3) locally verified; full revised arithmetic proof and final constants not certified.\par
Remaining | Sections 5–9 in full; external arithmetic and rational-point dependencies; complete Part I interface; journal-final comparison.\par
Adoption | No author contact or author adoption; focused AI acceptance complete, human adjudication pending.\par
The official arXiv record confirms v2. A bounded correction search found no directly relevant correction and does not establish absence of one. Findings are confined to the supplied author version.\par
\section*{Supporting checks}
The main assistant performed all scientific reasoning, static author-source inspection, and original-page inspection. Reviewer-owned exact-rational checks, matrix enumeration, reduction counts, source search, extraction/render logs, actual image ledger, and provenance records are saved. No author code was executed.\par
Backend equivalence and personal browser observation are not asserted. All writes stayed within this task’s staging directory.\par
The original compiler-absence probe preceded the authorized portable Tectonic follow-up. Its build and page inspections remain preserved as historical evidence. The current report is compiled from this exact TeX with Tectonic 0.17.0 in untrusted, cached-only mode. Exact input, dependency, log, PDF and page-image hashes, together with actual inspection of changed or reflowed pages, are recorded in the accompanying provenance. Original-review actions retain their original attribution; focused AI acceptance is recorded separately. Scientific adjudication remains pending, with no human review or author adoption claimed.\par
\section*{Consulted references}
9 A. Smith, supplied The distribution of \(\ell^\infty\)-Selmer groups in degree \(\ell\) twist families II, arXiv:2207.05143v2, PDF and TeX. https://arxiv.org/abs/2207.05143v2. A. Smith, supplied The distribution of \(\ell^\infty\)-Selmer groups in degree \(\ell\) twist families I, arXiv:2207.05674v2; separately reviewed as M25-0035. https://arxiv.org/abs/2207.05674v2.\par
\end{document}
